\section{Assembly of the count}\label{sec:assembly}

\begin{proof}[Proof of Theorem~\ref{thm:main}]
We first count the nonendpoint points with denominator coprime to \(3\)
in one shell \(T/2<q\le T\). Fix the parameters of
Lemma~\ref{lem:itinerary} and one nonempty label class. Choose \(X\)
uniform from it. Let \(E\) and the real traces \(W_i\) be as in
Lemma~\ref{lem:small-traces}. Since the data \((E,W_1,\ldots,W_J,X)\)
include \(X\) and are functions of \(X\), the logarithm of the size of the
label class is
\[
 H(E)+\sum_{i=1}^J H(W_i\mid E,W_1,\ldots,W_{i-1})
       +H(X\mid E,W_1,\ldots,W_J).
\]
The side information and the small steps together cost
\(O(\eta\log(2/\eta))\log T+O_\eta(1)\). We now bound the other
conditional entropies uniformly in the conditioning values that occur.

At a large step, use its dyadic upper bin endpoint \(U\). For each
\(u\) in the bin, the entire possible child trace lies above the set
\(A_u\) of Section~\ref{sec:boolean} with at most two lifts per point.
In this application use all eligible witnesses with the fixed past
and label; any further restrictions can only decrease the counts.
Discard empty fibers and put the nonempty ones into intervals of
width at most \(\Delta>0\) according to
\(\log |B_u|/\log U\). If the lower end of such an interval is \(\beta\),
then Proposition~\ref{prop:tariff} bounds the number of moduli by
\(T^{\xi+o(1)}U^{(1-\beta)/2+\sigma}\). By
\eqref{eq:bool-container}, each contributes at most
\(T^{o(1)}U^{\kappa(\beta+\Delta)+\kappa/(P_0-1)}\) points. Since
\(\kappa>1/2\), their total over all intervals, including the lifts, is
\[
 \ll_{\Delta} T^{\xi+o(1)}
      U^{\kappa+\sigma+\kappa\Delta+\kappa/(P_0-1)} .
\]
At the terminal step, for \(U\ge T^\eta\) use the same argument with
the bound \eqref{eq:bool-terminal}, and omit \(T^\xi\) and the
\(P_0\)-dependent power of \(U\). If the terminal upper bin endpoint
is less than \(T^\eta\), there are at most \(O(U^2)\) possibilities:
there are at most \(U\) moduli and at most \(U\) points per affine
grid. Each conditional entropy is at most the logarithm of the
corresponding support bound.

There are \(O(1/\eta)\) steps. The product of the true large increments
and the terminal modulus is at most \(q\), so the sum of the logarithms
of the corresponding bin upper endpoints is at most
\(\log T+O(1/\eta)\). The total extra loss outside the powers of the bin
endpoints is at most
\(O(\eta\log(2/\eta)+\eta+\xi/\eta)\log T+o(\log T)\), at fixed
parameters. To obtain any prescribed positive slack, first choose
\(\eta\) small and set \(\xi=\eta^3\). Then choose \(\sigma\) and
\(\Delta\) small. Follow the choice of parameters in the proof of
Proposition~\ref{prop:tariff}: fix the digit-entropy tolerances and
number of copies before the covering exponent; next use a fixed
finite covering menu and a fixed moment order for
Lemma~\ref{lem:lcm}. Finally choose \(P_0\) large enough for the
resulting finite number of inputs in Theorem~\ref{thm:native}, as well
as for the loss \(\kappa/(P_0-1)\). Only after these choices let \(T\)
grow. The stationary label count and the rotation losses are then
\(T^{o(1)}\); no claim of stationarity under conditioning is used.
This proves the bound with arbitrary positive slack for one shell.
Geometric summation gives the same bound for all purely periodic
points up to height \(T\), including the two endpoints at \(q=1\).

For a general reduced denominator \(q=3^j d\), with \(3\nmid d\),
deleting the first \(j\) ternary digits gives a point of reduced
denominator \(d\) and leaves at most \(2^j\) possible prefixes.
The integer tails \(0,1\) obey the same upper count. If the pure count
is \(O_\varepsilon(Y^{\kappa+\varepsilon})\) at height \(Y\ge1\),
the full count is bounded by a constant times
\[
 \sum_{0\le j\le\log_3 T} 2^j
                  (T/3^j)^{\kappa+\varepsilon}
 \le T^{\kappa+\varepsilon}\sum_{j\ge0}3^{-j\varepsilon}.
\]
This proves the theorem.
\end{proof}

The preperiod decomposition gives an exact identity. If \(P(Y)\)
counts the points up to reduced height \(Y\) with \(3\nmid q\), then
\begin{equation}\label{eq:preperiod}
 N_C(T)=P(T)+\sum_{j=1}^{\lfloor\log_3 T\rfloor}2^{j-1}P(T/3^j).
\end{equation}
In fact, write a pure tail as \(a/d=0.\overline B_3\), with
\(b\in\{0,2\}\) the last trit of \(B\). This includes the constant
tails \(0,1\). Modulo \(3\) one has \(a/d=-b\), by using the displayed
periodic denominator \(3^{|B|}-1\). Attach a prefix of length \(j\)
and integer value \(A\). The new numerator \(dA+a\) is coprime
to \(d\), and it is prime to \(3\) exactly when the last trit of
\(A\) differs from \(b\). Exactly half the prefixes thus give reduced
denominator \(3^j d\), and every point of this exact \(3\)-adic
valuation is recovered by stripping \(j\) trits.

In particular \(P(Y)\ll Y^\kappa\) would give the matching upper
bound \(N_C(T)\ll T^\kappa\log(2T)\). In fact, at \(T=3^m\) put
\(F_m=2^{-m}P(3^m)\). Identity~\eqref{eq:preperiod} reads
\(2^{-m}N_C(3^m)=F_m+\tfrac12\sum_{k<m}F_k\).
Consequently bounded averages of the nonnegative \(F_m\) are
equivalent to the matching logarithmic upper bound, by monotonicity
between triadic heights. The estimates above do not imply such
bounded averages.
